% 使用 article + xeCJK 替代 ctexart
% ctexart 会在 \begin{document} 时尝试加载系统字体集（SimSun 等），
% 在无系统字体或字体名不可解析的环境中导致 fontspec 失败。
% article + xeCJK 让字体设置完全由用户控制。
\documentclass[12pt,a4paper]{article}
\usepackage{xeCJK}

% === 字体设置（开源字体，便携免安装）===
% 通过 .otf/.ttf 扩展名触发 kpathsea 文件查找，无需 fontconfig
\setCJKmainfont{FandolSong-Regular.otf}[
  BoldFont=FandolSong-Bold.otf,
  ItalicFont=FandolKai-Regular.otf]
\setCJKsansfont{FandolHei-Regular.otf}[
  BoldFont=FandolHei-Bold.otf]
\setCJKmonofont{FandolKai-Regular.otf}
\setmainfont{texgyretermes-regular.otf}[
  BoldFont=texgyretermes-bold.otf,
  ItalicFont=texgyretermes-italic.otf,
  BoldItalicFont=texgyretermes-bolditalic.otf,
  Ligatures=TeXOff]
\setsansfont{texgyretermes-regular.otf}
\setmonofont{DejaVuSans.ttf}

% === 符号回退字体（FandolSong 不含 ★☆ 等特殊符号）===
\newfontfamily{\fallbacksymbols}{DejaVuSans.ttf}[Scale=MatchUppercase]

% === 宏包 ===
\usepackage{amsmath}
\usepackage{amssymb}
\usepackage[version=4]{mhchem}
\usepackage{graphicx}
\usepackage{paracol}
\usepackage{xcolor}
\usepackage{fancyhdr}
\usepackage{geometry}
\usepackage{tikz}

% === 红色圈号标记（TikZ 绘制，支持任意数字）===
\newcommand{\redcircled}[1]{%
  \tikz[baseline=(char.base)]{%
    \node[shape=circle,draw=red,inner sep=0.3pt,text=red,line width=0.4pt,font=\fontsize{7pt}{7pt}\selectfont] (char) {#1};%
  }%
}

% === 左栏标记：红色文字 + 圈号上标 ===
% 不包裹任何盒子，纯文本颜色，天然支持 CJK + 数学换行
\newcommand{\corrmark}[2]{%
  \textcolor{red}{#1}\textsuperscript{\textcolor{red}{\redcircled{#2}}}%
}

% === 页面设置 ===
\geometry{
  a4paper,
  left=1.5cm,
  right=1.5cm,
  top=2cm,
  bottom=2cm,
  columnsep=1cm
}

% === 内部辅助：计算 paracol 列宽 ===
\newlength{\colwidthinner}
\newcommand{\calccolwidth}{%
  \setlength{\colwidthinner}{\dimexpr\hsize-2\fboxsep-2\fboxrule\relax}%
}

% 右栏：修改建议框
\newcommand{\correctionbox}[1]{%
  \noindent\calccolwidth
  \colorbox{blue!8}{\parbox[t]{\colwidthinner}{#1}}%
}

% === Pandoc 兼容：非标准 LaTeX 命令 ===
\newcommand{\lt}{{<}}
\newcommand{\gt}{{>}}

% === 页眉页脚 ===
\pagestyle{fancy}
\fancyhf{}
\fancyhead[L]{校对报告}
\fancyhead[R]{\leftmark}
\fancyfoot[C]{\thepage}

\begin{document}

\title{校对报告}
\author{}
\date{}
\maketitle

{{CONTENT}}

\end{document}
